\begin{lemma}[A fixed exceptional span for a short arc]\label{lem:short-arc-exceptional-span}
For the short arc of length \(a\) in the cut comparison,
one may choose an exceptional set \(Z_c\) of target classes, independent of the input, with
\[
 \#Z_c\le(a+1)ma.
\] Outside its span the scaled lift has
exponent \(\sigma+E+\kappa_c\). The signed count is
\(r_c=a\) because each letter has an off-central window.
\end{lemma}

\begin{proof}
\label{par:short-exceptions}
On the marked target over the coefficient field let \(U\) be the
intersection of the transported \(\mathcal F_j^{u_j}\) over
positive levels. Let \(V_-\) be the sum of
\(\mathcal F_j^{-\eta}\) over negative levels. By the bound on
central multiplicities,
\[
 \operatorname{codim}U+\dim V_-\le ma .
\]
In the initial integral basis, the constants form the lattice
\(D_O\) with the transport of scalars to the terminal crystalline
space understood. Split the saturated sublattice \(D_O\cap U\)
off with complement \(K\), and project \(V_-\) along \(K\)
to \(U\). The intersection of the projected space with the
lattice of \(U\) is again saturated. Use adapted bases of these
two splittings.

Let \(e_j\) denote the adapted basis and let \(J\) index the complement \(K\) and this last intersection. Then
\[
 |J|\le ma,\qquad V_-\subseteq\operatorname{span}_{K_O}\{e_j:j\in J\}.
\]
Every other column belongs to \(U\).
The two changes of basis are integral and invertible, so they
give a Wach frame at \(c'\). Take the unique horizontal columns
with these initial values.

Apply the cut construction of Section~\ref{sec:inert-long-orbit}
to a crystalline extension representing a finite target class.
Its cut vector is the difference in that extension of a vector
in the incoming Frobenius span and a terminal Wach vector.
At each cyclotomic level the incoming span is regular; a pole
of order \(k>0\) therefore has coefficient in the extension's
\(\mathcal F^k\), and its image lies in the kernel of projection
to the trivial representation. Strictness identifies this
intersection with \(\mathcal F^k\) of the target.

For a polar term \(a_{-k}t^{-k}\) with \(k>0\), strictness therefore gives the two respective tests
\[
 a_{-k}\in V_-\quad(\text{negative letter}),\qquad k\le u_j+S\quad(\text{positive letter}).
\]
The statement for an entire polar part is obtained
coefficient by coefficient in the Laurent expansion, followed
by Chinese remaindering on the finite jets. At negative levels
all negative degrees thus disappear on indices outside \(J\).
Likewise \((\gamma-1)w\in XN_c\) imposes the Wach filtration
on the divided differences. Their projection outside \(J\)
vanishes at a negative level in degrees below \(u_j\):
for lower degrees the required filtration is contained in
\(\mathcal F_j^{-\eta}\), or equals it by the gap.

Write \(\mathcal J_c\) for the union of signed degree intervals
used in \(\kappa_c\).
Apply to these horizontal coordinates the integral operator
\[
 K(\gamma)=\left(\sum_{j=0}^{p^a-1}\gamma^j\right)
     \prod_{k\in\mathcal J_c\setminus\{0\}}
            (\gamma^{p^a}-\chi(\gamma)^{kp^a}) ,
\]
On a Laurent coefficient of
degree \(k\) at level \(t<a\), the principal action raised to
\(p^a\) is the scalar \(\chi(\gamma)^{kp^a}\). The sum
projects degree zero to a common value on the primitive roots
after tame invariants. Outside \(J\) this removes the excess
positive-level poles and the negative-level degrees from \(u_j\)
through \(u_j+S-1\).

At a negative level the divided difference condition on a nonconstant coefficient \(a_k\) below \(u_j\) reads
\[
 (\chi(\gamma)^{kp^a}-1)a_k=0,\qquad k\ne0.
\]
The scalar is nonzero, hence \(a_k=0\).
Thus the transformed scalar coefficients outside \(J\) satisfy
\eqref{eq:short-scalar}. The remaining constant at each negative
level is in the coefficient field, so a scalar power clears it.
The operators preserve the bounded integral denominator
\(\Lambda\); the improved pole orders descend to \(D\) by
division by the distinguished factors. Outside \(J\), the only poles allowed by \(D\) are positive with directions in \(U\); the negative zeros have the required orders. Write \(f=f_J+f_{J^c}\) for the support decomposition and \(\Delta\) for divided difference. The Wach test is linear, so
\[
 \Delta f_J=\Delta f-\Delta f_{J^c}
\]
passes as well.

For the latter part let \(M_J\) consist of the tuples with
denominator \(\Lambda\) supported on \(J\); let \(I_J\) be
the submodule whose divided differences pass the Wach test.
Let \(P_J\) consist of \(X\) times the tuples passing the Wach test
with bounded denominator, in \(M_J\). For numerators supported on \(J\),
a change by \(X\Lambda^2h\) induces the following change of tuple:
\[
 \frac{f+X\Lambda^2 h}{\Lambda}-\frac{f}{\Lambda}
       =X\Lambda h\in P_J .
\]
Weierstrass division reduces each numerator modulo \(X\Lambda^2\);
hence \(M_J/P_J\) is finite over \(O\).

The quotient \(M_J/P_J\) is also torsion-free.
Suppose \(0\ne c\in O\) and \(c(f/\Lambda)\in P_J\). Monic division by
\(X\) gives \(cf(0)=0\), hence \(f(0)=0\). The tuple
\(f/(X\Lambda)\) now has bounded denominator. Its rational filtration
tests hold since \(c f/(X\Lambda)\) satisfies them and the filtration
spaces are vector spaces over the coefficient field. Thus
\[
 f/\Lambda\in P_J.
\]


Let \(E_J\) be the image of \(I_J\) in \(M_J/P_J\).
After inverting \(p\), Chinese remainders for \(X\) and the
distinct cyclotomic factors of \(\Lambda\) split the numerator
jets at \(0\) and at the \(a\) levels. At \(0\) the quotient by
\(X\) leaves at most one entry per coordinate in \(J\).
At a cyclotomic level \(X\) is invertible, so the quotient is
by the filtered Laurent lattice.

In coefficient degree \(k\),
the principal power fixing the roots acts on that quotient by
\(\chi(\gamma)^{k p^a}\). A nonzero fixed class therefore needs
\(k=0\), since \(\chi(\gamma)\) has infinite order. In degree
zero the tame and principal root actions are transitive. Invariants
have at most one copy of each supported coordinate, and passing
to the quotient cannot increase this invariant dimension by
the characteristic-zero averaging projection. Thus
\[
    \dim_{\operatorname{Frac}(O)}(E_J[1/p])
          \le |J|+\sum_{t=0}^{a-1}|J|=(a+1)|J| .
\]
Since \(M_J/P_J\) is finite torsion-free over \(O\), this is
also a bound on the number of generators of \(E_J\).

Take an \(O\)-basis of the image and choose its representatives
in \(I_J\) once for the cut. Realize \(p^\sigma\) times each
by the horizontal comparison and take its cut Herr class. The
resulting set \(Z_c\) consists of finite classes, since a
product of cyclotomic factors carries the representative into
the Wach lattice. Using one representative per basis vector of \(E_J\) gives the uniform bound
\[
 |Z_c|\le\operatorname{rank}_O E_J\le (a+1)ma.
\]

\proofstep{Reduce every class against the same chosen representatives.}
For any input the exceptional coordinates are an integral
combination of these representatives modulo \(P_J\).
Terms of \(P_J\) represent elements of \(XN_c\) after the
same comparison and give zero under the cut class map.
The other coordinates lift at the common cost \(E\). To see the
exponent accounting, the inverse horizontal comparison reads
coordinates of \(p^\sigma w\). Realizing the split coordinates
requires another \(p^\sigma\), while on the outside part the
vector lifting relation already represents \(p^E\) times the
coordinates.

Since \(E\ge\sigma\), pad the realized split parts
by \(p^{E-\sigma}\). Their sum now matches
\[
 p^{\sigma+E} K(\gamma)w
\]
modulo \(XN_c\). Naturality and the vanishing on
\(XN_c\) follow from Frobenius contraction in the Herr complex,
as in the cut construction. Write \(K(T)-K(1)=(T-1)Q(T)\). Stability and the cut-vector condition give
\[
 K(\gamma)w-K(1)w=Q(\gamma)(\gamma-1)w\in XN_c.
\]
Thus the two vectors give the same cut class. In the scalar reading
\[
        p^{\sigma+E}K(1)w
             = u_c p^{\sigma+E+\kappa_c}w,\qquad u_c\in O^\times
\]
since \(v_p(K(1))=\kappa_c\). This gives the exponent
\(\sigma+E+\kappa_c\) with the same set \(Z_c\) for every
input class. For the present signed arcs every letter has a
positive or negative off-central window, so the signed level count
\(r_c=a\). \label{cor:inert-short-exceptions-1}
\end{proof}

